\section{Open Research Challenges}

\begin{frame}{}
    \LARGE Research Frontiers: \textbf{Open Research Challenges}

    \vspace{1em}
    \normalsize What is still unsolved, and how do you choose a problem worth your time?
\end{frame}

\begin{frame}{One Theme: Measurement Breaks When Optimised}
    \begin{center}
    \small
    \renewcommand{\arraystretch}{1.08}
    \setlength{\tabcolsep}{4pt}\begin{adjustbox}{max width=\linewidth}\begin{tabular}{lll}
        \toprule
        \textbf{Instrument} & \textbf{What happened once it became a target} & \textbf{Where} \\
        \midrule
        LLM judge & a lone ``:'' is accepted as a correct answer & Judges \\
        Rubric reward & the training score rises while the gold score falls & Judges \\
        Public leaderboard & private variants lift the rank, not the ability & Judges \\
        SAE reconstruction loss & better proxies, no better downstream use & Interpretability \\
        Speedup at batch size 1 & gains shrink or reverse under real load & Efficiency \\
        \bottomrule
    \end{tabular}\end{adjustbox}
    \end{center}
    \takeaway{The opportunities sit at the \textbf{measurement layer}: scoring that pays for abstention, judges hardened against attack, speedups reported at real batch sizes.}
    \src{Each row summarises a result shown with its source in the section named on the right.}
\end{frame}

\begin{frame}{Capabilities Are Jagged, and Benchmarks Miss It}
    \begin{itemize}\itemsep0.7em
        \item \textbf{Saturation.} Humanity's Last Exam gained 30 points in one year.
        \item \textbf{Jaggedness.} Models at Olympiad gold level read an analogue clock correctly about 50\% of the time. Humans: 90.1\%.
        \item \textbf{Broken items.} Invalid-question rates run from 2\% (MMLU Math) to 42\% (GSM8K).
        \item \textbf{Whose knowledge?} 28\% of MMLU questions need culturally specific knowledge.
    \end{itemize}
    \vspace{0.4em}
    \takeaway{A rising average hides uneven ability. Open problem: a benchmark that stays valid, culturally broad and unsaturated for more than a year.}
    \src{Stanford HAI, ``AI Index 2026,'' technical performance chapter (clock accuracy 50.6\% there, 50.1\% on its overview page). Singh et al., ``Global MMLU,'' \href{https://arxiv.org/abs/2412.03304}{arXiv:2412.03304}.}
\end{frame}

\begin{frame}{Can Scaling Continue? Compute, Data and Energy}
    \begin{columns}[T,onlytextwidth]
        \column{0.47\textwidth}
        \centering
        \small
        \renewcommand{\arraystretch}{1.2}
        \begin{adjustbox}{max width=\linewidth}\begin{tabular}{lc}
            \toprule
            \textbf{Constraint} & \textbf{Largest run, 2030} \\
            \midrule
            Power & $2\times10^{28}$ to $2\times10^{30}$ \\
            Chips & $10^{29}$ to $5\times10^{30}$ \\
            Data & $6\times10^{28}$ to $2\times10^{32}$ \\
            Latency & $3\times10^{30}$ to $10^{32}$ \\
            \midrule
            GPT-6 Astra (est.) & about $10^{27}$ \\
            \bottomrule
        \end{tabular}\end{adjustbox}

        \vspace{0.4em}
        {\scriptsize In FLOP. Epoch AI, 2024 analysis and 2026 estimate. Feasibility, not a forecast.}
        \column{0.51\textwidth}
        \begin{itemize}\itemsep0.45em
            \item Runs of about $2\times10^{29}$ FLOP look feasible by 2030. \textbf{Power binds first}, then chips.
            \item \textbf{Energy.} Data centres used about 415 TWh in 2024, projected at about 945 TWh by 2030.
            \item \textbf{Human text runs out.} About 300T usable tokens, fully used between 2026 and 2032.
            \item What is left is \textbf{experience}: interaction, rewards, continual learning.
        \end{itemize}
    \end{columns}
    \vspace{0.2em}
    \src{Sevilla et al., ``Can AI Scaling Continue Through 2030?'' and Villalobos et al., ``Will We Run Out of Data?'' (Epoch AI, 2024). IEA, ``Energy and AI.'' Silver, Sutton, ``Welcome to the Era of Experience.''}
\end{frame}

\begin{frame}{The Arabic Frontier: Standard Arabic Is Close. The Rest Is Open.}
    \begin{center}
    \small
    \renewcommand{\arraystretch}{1.18}
    \begin{adjustbox}{max width=\linewidth}\begin{tabular}{>{\raggedright\arraybackslash}p{3.9cm}>{\raggedright\arraybackslash}p{9.4cm}}
        \toprule
        \textbf{Evidence (date)} & \textbf{Finding} \\
        \midrule
        Open Arabic LLM Leaderboard v2 (2025) & v1 had machine-translated data and a silent scoring bug worth up to 20 points \\
        QuranicMMLU (Sep 2026) & 12 systems average 84\% on multiple choice, and 60\% on open-ended answers \\
        Mizan (Sep 2026) & each of 27 systems drops 14 to 18 points from Standard Arabic to Iraqi dialect \\
        \rowcolor{KAUSTcyan!12} Medical knowledge study (2026) & Arabic medical knowledge is \emph{present in intermediate layers} and fails to reach the output \\
        \bottomrule
    \end{tabular}\end{adjustbox}
    \end{center}
    \begin{itemize}\itemsep0.4em
        \item The highlighted row uses today's tools: a tuned lens and activation patching.
    \end{itemize}
    \src{El Filali et al., ``The Open Arabic LLM Leaderboard 2'' (2025). El Ghali et al., ``QuranicMMLU,'' \href{https://arxiv.org/abs/2609.22038}{arXiv:2609.22038}. Alseelawi, Aljumaily, ``Mizan,'' \href{https://arxiv.org/abs/2609.13980}{arXiv:2609.13980}. Abouzahir et al., Findings of EMNLP 2026, \href{https://arxiv.org/abs/2608.00207}{arXiv:2608.00207}.}
\end{frame}

\begin{frame}{What Is Missing? Leading Researchers Disagree}
    \begin{center}
    \small
    \renewcommand{\arraystretch}{1.2}
    \begin{adjustbox}{max width=\linewidth}\begin{tabular}{>{\raggedright\arraybackslash}p{2.3cm}>{\raggedright\arraybackslash}p{5.4cm}>{\raggedright\arraybackslash}p{5.0cm}}
        \toprule
        \textbf{Who (opinion)} & \textbf{What they say is missing} & \textbf{A test it suggests} \\
        \midrule
        Karpathy (2025) & continual learning. Outcome RL is ``sucking supervision through a straw'' & diversity loss over rounds of self-training \\
        Sutskever (2025) & generalisation, despite high benchmark scores & transfer from RL benchmark gains to held-out tasks \\
        Sutton, Silver & goals, and learning from experience & agents that learn from streams of reward \\
        LeCun & a world model that predicts representations, not tokens & replicate a claimed 16x data saving beyond one dataset \\
        Hassabis (2025) & consistency, and posing conjectures & invent a game as deep as Go \\
        \bottomrule
    \end{tabular}\end{adjustbox}
    \end{center}
    \src{Dwarkesh Podcast: Karpathy (2025-10-17), Sutskever (2025-11-25), Sutton (2025-09-26). Silver, Sutton, ``Welcome to the Era of Experience.'' Lex Fridman Podcast: LeCun (\#416), Hassabis (\#475). Huang, LeCun, Balestriero, \href{https://arxiv.org/abs/2602.22617}{arXiv:2602.22617}.}
\end{frame}
